\documentclass[11pt]{article}
\usepackage[margin=1in]{geometry}
\usepackage{booktabs}
\usepackage{amsmath}
\usepackage[T1]{fontenc}

\title{\vspace{-1.5em}AIST4010 Assignment 2: Antibiotic Resistance Gene (ARG) Classification\\[0.3em]
\large Experiment Report}
\author{Jiang Huanyun}
\date{November 4, 2026}

\begin{document}
\maketitle
\vspace{-2.5em}

\section*{1 What the task was}

15-class protein sequence classification from amino-acid FASTA: given a
protein, predict whether it is an antibiotic resistance gene (ARG) and, if
so, which of 14 antibiotic families it confers resistance to (aminoglycoside,
MLS, polymyxin, fosfomycin, trimethoprim, bacitracin, quinolone, multidrug,
chloramphenicol, tetracycline, rifampin, beta\_lactam, sulfonamide,
glycopeptide; label 14 = non-ARG). Train 21{,}209 sequences (10{,}067 ARG
with drug annotation, 11{,}142 non-ARG), val 4{,}387 (same skewed mix),
test 4{,}469. Metric is accuracy, and the test set is \emph{class-balanced}
($\sim$298 per class). I ended with a two-model ensemble of class-balanced
logistic-regression probes on ESM-2 650M and 3B embeddings: 98.10\% accuracy
on a balanced proxy of val (macro-F1 0.9753).

\section*{2 What I tried, in order}

(See Table~1 at the end of the document.)

\section*{3 Things that broke and how I fixed them}

\subsection*{3.1 Misreading the task as binary classification}

The training headers split ARG proteins by drug family
(\texttt{>id|FEATURES|src|drug|gene|mechanism|label}), and I initially read
the task as a binary ARG-vs-non-ARG problem, training on the trailing 0/1.
The first submission scored 0.94\%. The sample\_submission file shows 15
distinct labels (0--14). I then found the same assignment from Spring 2024
and another student's 2026 repo: the label space is 14 drug families plus
non-ARG, 15 classes in total, and this year's mapping matches the 2024
one. Lesson: before modeling, check the label vocabulary of the submission
template---it defines the problem, not the training-file format.

\subsection*{3.2 Validating on a skewed val fooled me twice}

Train and val share the same skewed mix (52\% non-ARG), but the test set is
balanced. Every model that learned the training prior scored great on val
(98.8--99.1\%) but produced heavily non-ARG-skewed test predictions. On a
balanced test that caps accuracy at $\sum_c \min(p_c, 298)/4469 \approx
38\%$ --- my second submission regressed. Two fixes: (i) train with
\emph{inverse-frequency class weights} so the model learns a uniform prior
that matches the test; (ii) evaluate on a \emph{balanced proxy} of val
(uniform 200 per class) and select models/weights by macro-F1 there, never
by raw val accuracy.

\subsection*{3.3 Fine-tuning hurt instead of helping}

LoRA fine-tuning of ESM-2 650M with a class-balanced cross-entropy reached
98.43\% raw val accuracy (macro-F1 0.954), and ensembling it with the
probe models pushed raw val accuracy to 99.13\% --- a number that looks
great but is distribution-mismatched. On the balanced proxy, the
fine-tuned model (97.47\%) was \emph{worse} than the plain balanced probe
(98.10\%), and adding it to the ensemble never helped. The fine-tuner
overfits the skewed training distribution; freezing ESM-2 and fitting a
linear probe regularizes better at this data scale.

\subsection*{3.4 Small bugs}

\begin{itemize}
\item \texttt{liblinear} doesn't support multiclass; the first 15-class
probe crashed, switched to \texttt{lbfgs}.
\item 3B checkpoint lacks tokenizer files (huggingface layout) --- copied
vocab/special tokens from the 650M checkpoint.
\item hf-mirror rate-limits (429) and huggingface.co direct TLS resets;
eventually downloaded the 3B shards with plain \texttt{curl -C -} resume.
\end{itemize}

\section*{4 What worked and what didn't}

\subsection*{4.1 Worked (roughly in order of contribution)}

\begin{enumerate}
\item ESM-2 650M mean-pool embeddings + LR (97.8\% balanced-proxy, the
single biggest jump over the k-mer baseline's 94.7\%);
\item class-balanced training (fixes the train/test prior mismatch, the
difference between 38\%-capped and 98\%);
\item balanced-proxy evaluation (the model selector that prevents choosing
skewed-prior models);
\item ensembling the 650M and 3B probes 1:1 (+0.3 on balanced proxy via
error decorrelation);
\item 3-mer TF-IDF + LR as a cheap diversity member and sanity baseline.
\end{enumerate}

\subsection*{4.2 Didn't work}

\begin{enumerate}
\item fine-tuning ESM-2 (LoRA) as a member of the ensemble (worse than the
probe, no ensemble gain);
\item ESM-2 3B probe alone (no gain over 650M; only helps in the
ensemble);
\item 4-gram features (no gain over 3-gram);
\item prior correction (multiplying class marginals to uniform) as a
post-hoc fix instead of class-weighted training.
\end{enumerate}

\subsection*{4.3 What I took away from it}

\begin{itemize}
\item The submission template's label distribution is part of the problem
specification: it reveals the test prior, and a model that ignores it can
look great locally and fail on the leaderboard.
\item For skewed-train/balanced-test setups, pick models by macro-F1 on a
balanced proxy of the validation set --- raw accuracy on the skewed val is
a biased selector.
\item Frozen pretrained embeddings + linear probe beat end-to-end
fine-tuning when the label distribution is skewed and the task is
near-linearly separable.
\item ``Strong and different'' still applies: the two probes share the same
architecture but different ESM-2 scales, and their errors decorrelate
enough to pay off.
\end{itemize}

\newpage
\begin{table}[h]
\centering
\small
\begin{tabular}{p{0.45cm} p{6.2cm} p{1.9cm} p{1.9cm} p{4.6cm}}
\toprule
\# & Method & Balanced-proxy acc & Worked? & Notes \\
\midrule
1 & 3-mer TF-IDF + LR, 15 classes & 94.73\% & baseline & fast sanity check; non-ARG vs multidrug confusion dominates \\
2 & ESM-2 650M embeddings + LR & 97.82\% & yes & frozen embeddings are already near-linearly separable \\
3 & ESM-2 650M + LoRA fine-tune (class-weighted) & 97.47\% & no & overfits skewed train prior; worse than the probe \\
4 & 4-gram features added & -- & no & no gain over 3-gram \\
5 & ESM-2 3B embeddings + LR & 98.17\% & slightly & $\sim$equal to 650M alone; gains only in the ensemble \\
6 & 650M + 3B probe ensemble (1:1) & 98.10\% & yes & error decorrelation between model scales \\
7 & probe + fine-tune + k-mer ensemble & 97.8--98.0\% & no & fine-tuner drags the ensemble down \\
8 & class-balanced training (inverse-freq weights) & -- & yes & the fix for the skewed-prior problem \\
9 & balanced-proxy model selection & -- & yes & the selector that exposed the val illusion \\
\bottomrule
\end{tabular}
\caption{Experiment log. Balanced-proxy acc: accuracy on a uniform 200-per-class subsample of val (the best available proxy for the balanced test set).}
\end{table}

\end{document}
